% !TEX TS-program = xelatex
% !TEX encoding = UTF-8 Unicode
% !Mode:: "TeX:UTF-8"
%
% 示例简历 —— 虚构人设"张三"，用于演示 my-resume.cls 的所有宏。
% 编译：xelatex example_zh.tex

\documentclass{my-resume}
\usepackage{my-resume-zh} % 中文支持（HarmonyOS Sans SC，自带字体，无需系统安装）
\usepackage{linespacing_fix} % 去掉 section 之间多余的空白
\usepackage{cite}

\begin{document}
\pagenumbering{gobble}

% \photoheader 用于中文版（带照片）；英文版建议用 \headerNoPhoto（见 example_en.tex）。
% photo_placeholder.jpg 只是演示用的灰色占位图，换成你自己的证件照即可（建议 2.6cm 宽的证件照比例）。
\photoheader{%
  \name{张三}
  \contactInfo{+86 138-0000-0000}{zhangsan@example.com}{现居城市：上海}{意向城市：上海/北京 \textbar\ 数据分析}
}{photo_placeholder.jpg}{2.6cm}
% 如果不想放照片：把 \photoheader{...}{photo_placeholder.jpg}{2.6cm} 换成 \headerNoPhoto{...}，
% 并删掉最后一个参数（照片路径与宽度）。

\section{个人总结}
仅 master（全经历版）需要这一段；岗位定制的 1 页版不写个人总结，把空间留给经历本身。写 2-4 句话：教育背景 + 1-2 项最有说服力的成果 + 想要的方向。

\section{教育背景}
\entry{某大学 (QS XX)}{2023.09 - 2025.07}{数据科学 硕士}{上海}
\begin{itemize}[parsep=0.2ex]
  \item \textbf{课程内容}：机器学习、统计推断、数据库系统、分布式计算
  \item \textbf{毕业设计}：题目 + 一句话说明方法/数据规模/核心发现
\end{itemize}
\entry{某大学 (QS XX)}{2019.09 - 2023.07}{计算机科学 本科}{北京}
\begin{itemize}[parsep=0.2ex]
  \item \textbf{相关课程}：数据结构与算法、概率论、线性代数、数据库原理
\end{itemize}

\section{实习经历}
\entry{某科技公司}{2024.07 - 2024.09}{数据分析实习生 \textperiodcentered\ 增长部}{上海}
\begin{itemize}[parsep=0.2ex]
  \item \textbf{加粗领起词 + 冒号}：一句话说明做了什么（用强动词，如"独立完成"/"主导"/"设计并实现"，避免"负责"）
  \item \textbf{量化结果}：只写能验证的真实数字；没有数字就保持定性描述，绝不编造
\end{itemize}

\section{科研与项目经历}
\entry{项目标题}{2024.01 - 2024.06}{}{上海}
\role{用到的工具/技术栈，逗号分隔}{}
\begin{itemize}[parsep=0.2ex]
  \item \textbf{方法/设计}：做了什么、为什么这么做
  \item \textbf{结果/验证}：量化的结果，或严谨的验证过程
  \item \textbf{影响/产出}：这个结果被用在哪、解决了什么问题
\end{itemize}

\section{技能/其他}
\entrySkillsSep
\begin{itemize}[parsep=0.2ex]
  \item \textbf{编程与工具}：Python，SQL，Git
  \item \textbf{数据分析}：pandas，统计建模，A/B 测试
  \item \textbf{语言}：英语（雅思 7.5），普通话（母语）
\end{itemize}

\end{document}
